% Generated draft: requires mathematical and source review.
\documentclass[11pt]{article}
\usepackage[margin=25mm]{geometry}
\usepackage{fontspec}
\setmainfont{Latin Modern Roman}
\usepackage{amsmath,amssymb,mathtools,unicode-math}
\setmathfont{Latin Modern Math}
\usepackage{enumitem,needspace,xurl,hyperref}
\setlength{\parindent}{0pt}
\setlength{\parskip}{5pt}
\setlength{\emergencystretch}{3em}
\newcommand{\N}{\mathbb{N}}
\newcommand{\Z}{\mathbb{Z}}
\newcommand{\Q}{\mathbb{Q}}
\newcommand{\R}{\mathbb{R}}
\newcommand{\C}{\mathbb{C}}
\newcommand{\F}{\mathbb{F}}
\newcommand{\A}{\mathbb{A}}
\newcommand{\PP}{\mathbb{P}}
\newcommand{\E}{\mathbb{E}}
\newcommand{\eps}{\varepsilon}
\newcommand{\dd}{\,\mathrm{d}}
\begin{document}
\section*{UnsolvedMath ID: 11000325}
\textbf{Source problem code:} catalogue entry
\subsection{Definitions and mathematical statement}
Let \(\mathcal H_{g,1}\) be the source's homology-cylinder group, whose precise surface marking and cobordism equivalence are defined in missing preceding text. This is not the handlebody subgroup of a mapping class group. Determine
\[
 (\mathcal H_{g,1})_{\mathrm{ab}}
   =\mathcal H_{g,1}/[\mathcal H_{g,1},\mathcal H_{g,1}],
 \qquad H_2(\mathcal H_{g,1};\Z).
\]
In particular, is the abelianization trivial, and is
\(\dim_\Q(H_2(\mathcal H_{g,1};\Z)\otimes\Q)=1\), with that rational direction detected by the intended signature class? Rank one does not assert that the full integral group is \(\Z\) or has no torsion; the signature normalization is not given.

The source conditionally expects nontrivial secondary homomorphisms
\[
 \widehat t_{2k+1}:\Theta^3_{\Z}\longrightarrow\Z
\]
associated to classes \(\widetilde t_{2k+1}\). The group \(\Theta^3_{\Z}\) is the intended integral-homology-\(3\)-sphere cobordism group, with connected sum; smooth/topological and equivalence conventions must be supplied. It compares these proposed invariants with the Casson invariant interpreted in [74] as a secondary invariant for the first MMM class \(e_1\). It says both constructions relate to \(H^2(\mathfrak h_{g,1}^{\Q})\), but distinguishes the reported nonzero \(e_1\) on \(\mathcal H_{g,1}\) from the expected vanishing of all \(\sigma^*(\widetilde t_{2k+1})\) there. Definitions of \(\sigma\), the classes, coefficients and index range are absent; the expected vanishing and resulting homomorphisms are not asserted as established.

It further reports Matumoto [67] and Galewsky--Stern [23] relating triangulability of every topological manifold in dimensions \(n\geq7\) to splitting of the missing map (21), and proposes studying the mod-\(2\) structure of the missing extension (20), motivated by Birman--Craggs [3] and Johnson [44] in the mapping-class case. Neither extension is supplied, so no splitting or nonsplitting result is invented. The final announced problem is also absent.
\subsection{Short English statement}
Compute the homology-cylinder group's abelianization and full integral second homology, testing perfection and the signature rank-one expectation; retain the conditional secondary-invariant and mod-two-extension program.
\subsection{Sources}
\begin{itemize}
\item \emph{Problems on Mapping Class Groups and Related Topics}, edited by Farb, Problem 11.4 and continuation, supplied PDF page 379. \url{https://www.math.uchicago.edu/~farb/papers/mcgbook.pdf}. Recorded access 2026-07-29; extraction pdf-font-structure.
\item [74], Casson/secondary MMM interpretation; Matumoto [67]; Galewsky and Stern [23] (source spelling); Birman and Craggs [3]; Johnson [44]. Full bibliography, extensions (20), (21) and the class constructions are absent.
\item The supplied summary mentions Johnson, Mess and Hain as Torelli analogues, Garoufalidis--Levine tree/finite-type work and Sakasai's cylinder/Magnus/Johnson work without complete citations; it reports no completed determination or secondary-invariant construction.
\item Local UnsolvedMath registry, \texttt{data/unsolvedmath/problems.json}, permanent ID \(11000325\), exact source code \texttt{catalogue entry}; supplied OPEN-TRIAGE status.
\end{itemize}
\end{document}
